\section{Erweitertes Unpacking und Funktionen}
\label{sec:functions_unpacking}

\subsection{Tuple Packing und Unpacking}
Python ermöglicht das Gruppieren und Verteilen von Objekten über Mechanismen zum Packen und Entpacken.

\begin{itemize}
    \item \textbf{Packing:} Mehrere einzelne Objekte werden zu einem einzigen Tupel gruppiert.
    \item \textbf{Unpacking:} Elemente in einem iterierbaren Objekt werden einzelnen Variablen zugewiesen. Dafür muss die Anzahl der Variablen der Anzahl der Elemente entsprechen, sofern kein erweitertes Entpacken verwendet wird.
\end{itemize}

\begin{lstlisting}[language=Python]
# Packing
employee = 'Peter', 'Muller' # Ergibt ('Peter', 'Muller')

# Unpacking
first, last = employee

# Mehrfachzuweisung (kombiniertes Packing/Unpacking)
x, y, z = 5, 10, 15
\end{lstlisting}

\subsubsection{Erweitertes Iterable-Unpacking}
Das Entpacken funktioniert mit jedem iterierbaren Objekt (Strings, Listen, Mengen) und unterstützt bestimmte Operatoren, um variable Längen oder unerwünschte Werte zu behandeln.

\begin{itemize}
    \item \textbf{Rest-Operator (*):} Erfasst die verbleibenden Elemente in einer Liste.
    \item \textbf{Platzhalter (\_):} Wird konventionell verwendet, um bestimmte Werte beim Entpacken zu ignorieren.
\end{itemize}

\begin{lstlisting}[language=Python]
# Unpacking mit Rest
first, last, *address, birth_date = ("Peter", "Muller", "Street", 10, "14.04.1977")
# address wird ["Street", 10]

# Werte ignorieren
first, last, _, birth_date = employee_data

# Entpacken in verschachtelten Strukturen
x = [1, 2, (3, [4, 5])]
a, b, (c, (d, e)) = x
\end{lstlisting}

\subsubsection{Iterable-Expansion}
Die Operatoren \texttt{*} und \texttt{**} ermöglichen das Entpacken iterierbarer Objekte in anderen Collection-Literalen.

\begin{lstlisting}[language=Python]
list1 = [1, 2]
list2 = [*list1, 3, 4] # [1, 2, 3, 4]

dict1 = {"A": 1, "B": 2}
dict2 = {"C": 3, **dict1} # {'C': 3, 'A': 1, 'B': 2}
\end{lstlisting}

\subsection{Funktionen}
Funktionen werden mit dem Schlüsselwort \texttt{def} definiert. Sie legen einen lokalen Gültigkeitsbereich fest und können Objekte an den Aufrufer zurückgeben.

\begin{itemize}
    \item \textbf{Parameter:} Platzhalter, die in der Funktionssignatur definiert sind.
    \item \textbf{Argumente:} Tatsächliche Werte, die beim Aufruf übergeben werden.
    \item \textbf{Rückgabewert:} Funktionen geben standardmäßig \texttt{None} zurück. Mehrere Werte können als Tupel zurückgegeben werden.
\end{itemize}

\subsubsection{Argumenttypen und Einschränkungen}
\begin{itemize}
    \item \textbf{Positional vs. Keyword:} Positionsargumente müssen vor Keyword-Argumenten in einem Aufruf stehen.
    \item \textbf{Standardparameter:} Parameter mit Standardwerten müssen auf diejenigen ohne Standardwert folgen.
    \item \textbf{Variadische Argumente (*args, **kwargs):} \texttt{*args} sammelt zusätzliche positionsbezogene Argumente in ein 
    \crd{\textbf{tuple}}; \texttt{**kwargs} sammelt zusätzliche Keyword-Argumente in ein Dictionary.
\end{itemize}

\subsubsection{Argument-Erzwingung}
Entwickler können bestimmte Aufrufstile erzwingen, indem sie \texttt{/} und \texttt{*} verwenden.

\begin{lstlisting}[language=Python]
def constrained_func(pos_only, /, pos_or_kw, *, kw_only):
    pass

# / : Alle Argumente vor müssen positionsspezifisch sein.
# * : Alle Argumente nach müssen keyword-only sein.
\end{lstlisting}

\subsection{Docstrings}
Dokumentationsstrings sollten der Funktionsdefinition folgen, um ihr Verhalten zu beschreiben. Sie können über \texttt{help(function\_name)} oder das \texttt{\_\_doc\_\_}-Attribut abgerufen werden.

\begin{lstlisting}[language=Python]
def add(a, b):
    """Gibt die Summe von a und b zurück."""
    return a + b
\end{lstlisting}

\subsection{Lambda-Funktionen}
Lambdas sind anonyme Funktionen mit einem einzigen Ausdruck. Sie werden hauptsächlich als "key"-Funktionen zum Sortieren oder Filtern verwendet.

\textbf{Syntax:} \texttt{lambda parameters: expression}

\begin{lstlisting}[language=Python]
# Sortieren einer Liste von Strings nach dem numerischen Teil
articles = ['art1', 'art20', 'art3']
sorted_list = sorted(articles, key=lambda x: int(x[3:]))
\end{lstlisting}